% GENERATED FILE. Edit canonical lessons, then run npm run generate:books.
% canonical-source-hash: fnv1a64:7e7840c05c910406
% canonical-lessons: KA-C326-kaneyagu, KA-C326-patte-haccu, KA-C326-kuruhu, KA-C326-apaharisu, KA-C326-otteyalu, KA-R326-review-319-326

\chapter{A Child Goes Missing}
\label{ch:ka-a-child-goes-missing}
\hlchaptermodality{326}

\begin{chapteropening}
\textbf{By the end of this chapter:} \emph{I can say to go missing, to track down, a trace, to kidnap and a hostage.}

Complete the last lesson of chapter 326: I can say to go missing, to track down, a trace, to kidnap and a hostage.
\end{chapteropening}

\section[\texorpdfstring{\kn{ಕಾಣೆಯಾಗು} (kāṇeyāgu)}{ಕಾಣೆಯಾಗು (kāṇeyāgu)}]{\kn{ಕಾಣೆಯಾಗು} (kāṇeyāgu) — to go missing}

\label{lesson:KA-C326-kaneyagu}



\begin{quote}
Before the new one: say the Kannada for to bless, then the Kannada for a wedding party.
\end{quote}

\subsection*{You'll want to know: \kn{ಕಾಣೆಯಾಗು}}
\textbf{\kn{ಕಾಣೆಯಾಗು}} — \emph{kāṇeyāgu} — \textquotedblleft{}to go missing\textquotedblright{}.

\begin{quote}
\kn{ಒಂದು} \kn{ವೇಳೆ} \kn{ಮಗು} \kn{ಕಾಣೆಯಾದರೆ} \kn{ಅಪ್ಪ} \kn{ಅಮ್ಮ} \kn{ತಕ್ಷಣ} \kn{ದೂರು} \kn{ಕೊಡುತ್ತಾರೆ}.
\end{quote}

\begin{quote}
\emph{ondu vēḷe magu kāṇeyādare appa amma takṣaṇa dūru koḍuttāre.}
\end{quote}

\begin{quote}
If a child goes missing, the parents report it at once.
\end{quote}

Past (it …): \kn{ಕಾಣೆಯಾಯಿತು}, \emph{kāṇeyāyitu}.

\kn{ಕಾಣೆ}, not to be seen, and \kn{ಆಗು}, to become.

\begin{cousinweb}
\textbf{\kn{ಕಾಣೆಯಾಗು}} (\emph{kāṇeyāgu}) is a native Dravidian compound: \kn{ಕಾಣೆ}, \textquotedblleft{}not seen\textquotedblright{}, from \emph{kāṇ} \textquotedblleft{}to see\textquotedblright{} (DEDR 1443), and \emph{āgu} \textquotedblleft{}to become\textquotedblright{} (DEDR 333).
\end{cousinweb}

\begin{grammarlens}[title={What you've built}]
One more word for telling what happened.
\end{grammarlens}

\subsection*{Your turn}
\begin{itemize}
\raggedright
  \item \emph{Say it:} \emph{kāṇeyāgu}
  \item \emph{Say it:} \emph{kāṇeyāgu}, once more
  \item \emph{Say it:} \emph{ondu vēḷe magu kāṇeyādare appa amma takṣaṇa dūru koḍuttāre.} — If a child goes missing, the parents report it at once.
  \item \emph{Recall:} say the Kannada for to bless, then the Kannada for a wedding party, then say \emph{kāṇeyāgu} again
\end{itemize}

\begin{tcolorbox}[breakable,colback=teal!4,colframe=teal!35!black,title={Before you move on}]
What does \kn{ಕಾಣೆಯಾಗು} mean? (\textquotedblleft{}To go missing\textquotedblright{}.) Say it once more.
\end{tcolorbox}
\section[\texorpdfstring{\kn{ಪತ್ತೆ} \kn{ಹಚ್ಚು} (patte haccu)}{ಪತ್ತೆ ಹಚ್ಚು (patte haccu)}]{\kn{ಪತ್ತೆ} \kn{ಹಚ್ಚು} (patte haccu) — to track down; to detect}

\label{lesson:KA-C326-patte-haccu}



\begin{quote}
Before the new one: say the Kannada for a wedding party, then the Kannada for to go missing.
\end{quote}

\subsection*{You'll want to know: \kn{ಪತ್ತೆ} \kn{ಹಚ್ಚು}}
\textbf{\kn{ಪತ್ತೆ} \kn{ಹಚ್ಚು}} — \emph{patte haccu} — \textquotedblleft{}to track down; to detect\textquotedblright{}.

\begin{quote}
\kn{ಪೊಲೀಸರು} \kn{ಎರಡು} \kn{ದಿನದಲ್ಲಿ} \kn{ಮಗುವನ್ನು} \kn{ಪತ್ತೆ} \kn{ಹಚ್ಚಿದರು}.
\end{quote}

\begin{quote}
\emph{polīsaru eraḍu dinadalli maguvannu patte haccidaru.}
\end{quote}

\begin{quote}
The police tracked down the child in two days.
\end{quote}

Past (I …): \kn{ಪತ್ತೆ} \kn{ಹಚ್ಚಿದೆ}, \emph{patte haccide}.

\kn{ಪತ್ತೆ} is a trace; \kn{ಪತ್ತೆಯಾಗು} means to be found.

\begin{cousinweb}
\textbf{\kn{ಪತ್ತೆ} \kn{ಹಚ್ಚು}} (\emph{patte haccu}) is a Kannada phrase verb built from \kn{ಪತ್ತೆ}, \textquotedblleft{}a trace, a whereabouts\textquotedblright{}, and \kn{ಹಚ್ಚು}, \textquotedblleft{}to apply, to light\textquotedblright{}. The origin of \kn{ಪತ್ತೆ} is not settled.
\end{cousinweb}

\begin{grammarlens}[title={What you've built}]
One more word for telling what happened.
\end{grammarlens}

\subsection*{Your turn}
\begin{itemize}
\raggedright
  \item \emph{Say it:} \emph{patte haccu}
  \item \emph{Say it:} \emph{patte haccu}, once more
  \item \emph{Say it:} \emph{polīsaru eraḍu dinadalli maguvannu patte haccidaru.} — The police tracked down the child in two days.
  \item \emph{Recall:} say the Kannada for a wedding party, then the Kannada for to go missing, then say \emph{patte haccu} again
\end{itemize}

\begin{tcolorbox}[breakable,colback=teal!4,colframe=teal!35!black,title={Before you move on}]
What does \kn{ಪತ್ತೆ} \kn{ಹಚ್ಚು} mean? (\textquotedblleft{}To track down; to detect\textquotedblright{}.) Say it once more.
\end{tcolorbox}
\section[\texorpdfstring{\kn{ಕುರುಹು} (kuruhu)}{ಕುರುಹು (kuruhu)}]{\kn{ಕುರುಹು} (kuruhu) — a trace; a clue}

\label{lesson:KA-C326-kuruhu}



\begin{quote}
Before the new one: say the Kannada for to go missing, then the Kannada for to track down.
\end{quote}

\subsection*{You'll want to know: \kn{ಕುರುಹು}}
\textbf{\kn{ಕುರುಹು}} — \emph{kuruhu} — \textquotedblleft{}a trace; a clue\textquotedblright{}.

\begin{quote}
\kn{ಕಳ್ಳ} \kn{ಯಾವ} \kn{ಕುರುಹೂ} \kn{ಬಿಡಲಿಲ್ಲ}.
\end{quote}

\begin{quote}
\emph{kaḷḷa yāva kuruhū biḍalilla.}
\end{quote}

\begin{quote}
The thief left no trace.
\end{quote}

\kn{ಕುರುಹು} \kn{ಬಿಡು}, to leave a trace; \kn{ಕುರುಹು} \kn{ಸಿಗು}, to find a clue.

\begin{cousinweb}
\textbf{\kn{ಕುರುಹು}} (\emph{kuruhu}) is native Dravidian: DEDR 1847 lists Kannada \emph{kuṟuhu}, \textquotedblleft{}a mark, a sign, a token\textquotedblright{}, beside \emph{kuṟi}, \textquotedblleft{}a mark, an aim\textquotedblright{}. The verb \kn{ಗುರುತಿಸು} belongs to the same root.
\end{cousinweb}

\begin{grammarlens}[title={What you've built}]
One more word for telling what happened.
\end{grammarlens}

\subsection*{Your turn}
\begin{itemize}
\raggedright
  \item \emph{Say it:} \emph{kuruhu}
  \item \emph{Say it:} \emph{kuruhu}, once more
  \item \emph{Say it:} \emph{kaḷḷa yāva kuruhū biḍalilla.} — The thief left no trace.
  \item \emph{Recall:} say the Kannada for to go missing, then the Kannada for to track down, then say \emph{kuruhu} again
\end{itemize}

\begin{tcolorbox}[breakable,colback=teal!4,colframe=teal!35!black,title={Before you move on}]
What does \kn{ಕುರುಹು} mean? (\textquotedblleft{}A trace; a clue\textquotedblright{}.) Say it once more.
\end{tcolorbox}
\section[\texorpdfstring{\kn{ಅಪಹರಿಸು} (apaharisu)}{ಅಪಹರಿಸು (apaharisu)}]{\kn{ಅಪಹರಿಸು} (apaharisu) — to kidnap; to abduct}

\label{lesson:KA-C326-apaharisu}



\begin{quote}
Before the new one: say the Kannada for to track down, then the Kannada for a trace.
\end{quote}

\subsection*{You'll want to know: \kn{ಅಪಹರಿಸು}}
\textbf{\kn{ಅಪಹರಿಸು}} — \emph{apaharisu} — \textquotedblleft{}to kidnap; to abduct\textquotedblright{}.

\begin{quote}
\kn{ಇಬ್ಬರು} \kn{ಆ} \kn{ಹುಡುಗನನ್ನು} \kn{ಅಪಹರಿಸಿದರು}.
\end{quote}

\begin{quote}
\emph{ibbaru ā huḍuganannu apaharisidaru.}
\end{quote}

\begin{quote}
Two people kidnapped that boy.
\end{quote}

Past (I …): \kn{ಅಪಹರಿಸಿದೆ}, \emph{apahariside}.

The person taken takes -\kn{ಅನ್ನು}. The noun is \kn{ಅಪಹರಣ}, a kidnapping.

\begin{cousinweb}
\textbf{\kn{ಅಪಹರಿಸು}} (\emph{apaharisu}) is built on Sanskrit \emph{apaharaṇa}, \textquotedblleft{}taking away, carrying off, stealing\textquotedblright{} (Monier-Williams), with the Kannada verb ending -\emph{isu}.
\end{cousinweb}

\begin{grammarlens}[title={What you've built}]
One more word for telling what happened.
\end{grammarlens}

\subsection*{Your turn}
\begin{itemize}
\raggedright
  \item \emph{Say it:} \emph{apaharisu}
  \item \emph{Say it:} \emph{apaharisu}, once more
  \item \emph{Say it:} \emph{ibbaru ā huḍuganannu apaharisidaru.} — Two people kidnapped that boy.
  \item \emph{Recall:} say the Kannada for to track down, then the Kannada for a trace, then say \emph{apaharisu} again
\end{itemize}

\begin{tcolorbox}[breakable,colback=teal!4,colframe=teal!35!black,title={Before you move on}]
What does \kn{ಅಪಹರಿಸು} mean? (\textquotedblleft{}To kidnap; to abduct\textquotedblright{}.) Say it once more.
\end{tcolorbox}
\section[\texorpdfstring{\kn{ಒತ್ತೆಯಾಳು} (otteyāḷu)}{ಒತ್ತೆಯಾಳು (otteyāḷu)}]{\kn{ಒತ್ತೆಯಾಳು} (otteyāḷu) — a hostage}

\label{lesson:KA-C326-otteyalu}



\begin{quote}
Before the new one: say the Kannada for a trace, then the Kannada for to kidnap.
\end{quote}

\subsection*{You'll want to know: \kn{ಒತ್ತೆಯಾಳು}}
\textbf{\kn{ಒತ್ತೆಯಾಳು}} — \emph{otteyāḷu} — \textquotedblleft{}a hostage\textquotedblright{}.

\begin{quote}
\kn{ಕಳ್ಳರು} \kn{ಅಂಗಡಿಯವನನ್ನು} \kn{ಒತ್ತೆಯಾಳಾಗಿ} \kn{ಇಟ್ಟುಕೊಂಡರು}.
\end{quote}

\begin{quote}
\emph{kaḷḷaru aṅgaḍiyavanannu otteyāḷāgi iṭṭukoṇḍaru.}
\end{quote}

\begin{quote}
The robbers kept the shopkeeper as a hostage.
\end{quote}

Plural: \kn{ಒತ್ತೆಯಾಳುಗಳು}, \emph{otteyāḷugaḷu}.

\kn{ಒತ್ತೆ}, a pledge, and \kn{ಆಳು}, a person: a person held as a pledge.

\begin{cousinweb}
\textbf{\kn{ಒತ್ತೆಯಾಳು}} (\emph{otteyāḷu}) is a native Dravidian compound: \emph{otte}, \textquotedblleft{}a pledge, a pawn\textquotedblright{} (DEDR 1020), and \emph{āḷ}, \textquotedblleft{}a person, a servant\textquotedblright{} (DEDR 399).
\end{cousinweb}

\begin{grammarlens}[title={What you've built}]
One more word for telling what happened.
\end{grammarlens}

\subsection*{Your turn}
\begin{itemize}
\raggedright
  \item \emph{Say it:} \emph{otteyāḷu}
  \item \emph{Say it:} \emph{otteyāḷu}, once more
  \item \emph{Say it:} \emph{kaḷḷaru aṅgaḍiyavanannu otteyāḷāgi iṭṭukoṇḍaru.} — The robbers kept the shopkeeper as a hostage.
  \item \emph{Recall:} say the Kannada for a trace, then the Kannada for to kidnap, then say \emph{otteyāḷu} again
\end{itemize}

\begin{tcolorbox}[breakable,colback=teal!4,colframe=teal!35!black,title={Before you move on}]
What does \kn{ಒತ್ತೆಯಾಳು} mean? (\textquotedblleft{}A hostage\textquotedblright{}.) Say it once more.
\end{tcolorbox}
\section[(dialogue)]{Telling what happened}

\label{lesson:KA-R326-review-319-326}



\begin{quote}
Say the Kannada for a war and a soldier.
\end{quote}

\subsection*{Your turn}
Say these aloud:

\begin{itemize}
\raggedright
  \item a war, a soldier, to invade, a thief, to run away
  \item a reporter, to occur, a broadcast, to broadcast, a headline
  \item a court of law, a judge, to accuse, a verdict, a punishment
  \item a king, a queen, to reign, a throne, a crown
  \item a crime, to arrest, a police station, an investigation, to suspect
  \item a firefighter, to catch fire, presence of mind, to extinguish, to rescue
  \item a marriage, a wedding reception, a dowry, to bless, a wedding party
  \item to go missing, to track down, a trace, to kidnap, a hostage
\end{itemize}

\begin{tcolorbox}[breakable,colback=teal!4,colframe=teal!35!black,title={Before you move on}]
A war? (\textbf{\kn{ಯುದ್ಧ}}, \emph{yuddha}.) A soldier? (\textbf{\kn{ಸೈನಿಕ}}, \emph{sainika}.) To invade? (\textbf{\kn{ಆಕ್ರಮಿಸು}}, \emph{ākramisu}.) A thief? (\textbf{\kn{ಕಳ್ಳ}}, \emph{kaḷḷa}.) To run away? (\textbf{\kn{ಓಡಿಹೋಗು}}, \emph{ōḍihōgu}.) A reporter? (\textbf{\kn{ವರದಿಗಾರ}}, \emph{varadigāra}.) To occur? (\textbf{\kn{ಸಂಭವಿಸು}}, \emph{sambhavisu}.) A broadcast? (\textbf{\kn{ಪ್ರಸಾರ}}, \emph{prasāra}.) To broadcast? (\textbf{\kn{ಬಿತ್ತರಿಸು}}, \emph{bittarisu}.) A headline? (\textbf{\kn{ಶೀರ್ಷಿಕೆ}}, \emph{śīrṣike}.) A court of law? (\textbf{\kn{ನ್ಯಾಯಾಲಯ}}, \emph{nyāyālaya}.) A judge? (\textbf{\kn{ನ್ಯಾಯಾಧೀಶ}}, \emph{nyāyādhīśa}.) To accuse? (\textbf{\kn{ಆರೋಪಿಸು}}, \emph{ārōpisu}.) A verdict? (\textbf{\kn{ತೀರ್ಪು}}, \emph{tīrpu}.) A punishment? (\textbf{\kn{ಶಿಕ್ಷೆ}}, \emph{śikṣe}.) A king? (\textbf{\kn{ರಾಜ}}, \emph{rāja}.) A queen? (\textbf{\kn{ರಾಣಿ}}, \emph{rāṇi}.) To reign? (\textbf{\kn{ಆಳು}}, \emph{āḷu}.) A throne? (\textbf{\kn{ಸಿಂಹಾಸನ}}, \emph{siṃhāsana}.) A crown? (\textbf{\kn{ಕಿರೀಟ}}, \emph{kirīṭa}.) A crime? (\textbf{\kn{ಅಪರಾಧ}}, \emph{aparādha}.) To arrest? (\textbf{\kn{ಬಂಧಿಸು}}, \emph{bandhisu}.) A police station? (\textbf{\kn{ಠಾಣೆ}}, \emph{ṭhāṇe}.) An investigation? (\textbf{\kn{ತನಿಖೆ}}, \emph{tanikhe}.) To suspect? (\textbf{\kn{ಶಂಕಿಸು}}, \emph{śaṅkisu}.) A firefighter? (\textbf{\kn{ಅಗ್ನಿಶಾಮಕ}}, \emph{agniśāmaka}.) To catch fire? (\textbf{\kn{ಹೊತ್ತಿಕೊ}}, \emph{hottiko}.) Presence of mind? (\textbf{\kn{ಸಮಯಪ್ರಜ್ಞೆ}}, \emph{samayaprajñe}.) To extinguish? (\textbf{\kn{ನಂದಿಸು}}, \emph{nandisu}.) To rescue? (\textbf{\kn{ರಕ್ಷಿಸು}}, \emph{rakṣisu}.) A marriage? (\textbf{\kn{ವಿವಾಹ}}, \emph{vivāha}.) A wedding reception? (\textbf{\kn{ಆರತಕ್ಷತೆ}}, \emph{āratakṣate}.) A dowry? (\textbf{\kn{ವರದಕ್ಷಿಣೆ}}, \emph{varadakṣiṇe}.) To bless? (\textbf{\kn{ಹರಸು}}, \emph{harasu}.) A wedding party? (\textbf{\kn{ದಿಬ್ಬಣ}}, \emph{dibbaṇa}.) To go missing? (\textbf{\kn{ಕಾಣೆಯಾಗು}}, \emph{kāṇeyāgu}.) To track down? (\textbf{\kn{ಪತ್ತೆ} \kn{ಹಚ್ಚು}}, \emph{patte haccu}.) A trace? (\textbf{\kn{ಕುರುಹು}}, \emph{kuruhu}.) To kidnap? (\textbf{\kn{ಅಪಹರಿಸು}}, \emph{apaharisu}.) A hostage? (\textbf{\kn{ಒತ್ತೆಯಾಳು}}, \emph{otteyāḷu}.)
\end{tcolorbox}
